%

\section{Zastosowanie: twierdzenie Urquharta} % lang-pl
\section{Applicazione: teorema di Urquhart} % lang-it

\begin{theorem}
[Urquharta] % lang-pl
[di Urquhart] % lang-it
\label{theorem_urquhart}%
\index{twierdzenie!Urquharta}% % lang-pl
\index{teorema!di Urquhart}% % lang-it
    Przy oznaczeniach jak na rysunku: % lang-pl
    Con le notazioni della figura: % lang-it
    \begin{center}
\begin{comment}
    \begin{tikzpicture}[scale=.4]
        \tkzDefPoint(0, 0){E}
        \tkzDefPoint(267:5){E1}
        \tkzDefPoint(250:5){E2}
        \tkzDefPoint(165:5){E3}
        \tkzDefPoint(135:5){E4}

        \tkzDefLine[tangent at=E1](E) \tkzGetPoint{tan1}
        \tkzDefLine[tangent at=E2](E) \tkzGetPoint{tan2}
        \tkzDefLine[tangent at=E3](E) \tkzGetPoint{tan3}
        \tkzDefLine[tangent at=E4](E) \tkzGetPoint{tan4}

        \tkzInterLL(E4,tan4)(E1,tan1) \tkzGetPoint{S1}
        \tkzInterLL(E4,tan4)(E2,tan2) \tkzGetPoint{S2}
        \tkzInterLL(E3,tan3)(E1,tan1) \tkzGetPoint{S3}
        \tkzInterLL(E3,tan3)(E2,tan2) \tkzGetPoint{S4}

        \tkzInterLL(S3,S4)(S1,S2) \tkzGetPoint{Gorny}
        \tkzInterLL(S1,S3)(S2,S4) \tkzGetPoint{Dolny}

        \tkzLabelPoint[below](S1){$A$}
        \tkzLabelPoint[above left](S2){$D$}
        \tkzLabelPoint(S3){$B$}
        \tkzLabelPoint[above right](S4){$C$}
        \tkzLabelPoint(Dolny){$E$}
        \tkzLabelPoint[right](Gorny){$F$}

        \tkzDrawSegments[line width=0.2mm](S1,Dolny S1,Gorny S3,Gorny S2,Dolny)
        \tkzDrawPoints[size=3,color=black,fill=black!50](S1,S2,S3,S4,Dolny,Gorny)
        % konstrukcja z https://en.wikipedia.org/wiki/Ex-tangential_quadrilateral
\end{tikzpicture}
\end{comment}
    \end{center}
    niech $|AB| + |BC| = |AD| + |CD|$. % lang-pl
    Wtedy $|AE| + |CE| = |AF| + |CF|$. % lang-pl
       sia $|AB| + |BC| = |AD| + |CD|$. % lang-it
    Allora $|AE| + |CE| = |AF| + |CF|$. % lang-it
\end{theorem}

Twierdzenie \ref{theorem_urquhart} nie będzie zbyt popularne, napiszą o nim Bogdańska, Neugebauer \cite[s. 97]{neugebauer_2018}, nieznany autor w $\Delta_{84}^{11}$ i pewnie ktoś jeszcze. % lang-pl
Nie do końca wiadomo, kto udowodni je pierwszy: być może będzie to de Morgan w 1841 roku; teza jest też przypadkiem granicznym innego wyniku Chaslesa, znanego około roku 1860. % lang-pl
Wreszcie Dan Pedoe \cite{pedoe_1976} przypisze to twierdzenie Malcolmowi Livingstonowi Urquhartowi\footnote{Matematyk australijski, będzie żyć w latach 1902-1966 i nie opublikuje żadnej pracy.}. % lang-pl

Il teorema \ref{theorem_urquhart} non diventerà particolarmente popolare; ne scriveranno Bogdańska e Neugebauer \cite[s. 97]{neugebauer_2018}, un autore anonimo in $\Delta_{84}^{11}$ e probabilmente anche qualcun altro. % lang-it
Non è del tutto chiaro chi ne fornirà per primo una dimostrazione: forse de Morgan nel 1841; la tesi è anche un caso limite di un altro risultato di Chasles, noto intorno al 1860. % lang-it
Infine Dan Pedoe \cite{pedoe_1976} attribuirà questo teorema a Malcolm Livingston Urquhart\footnote{Matematico australiano, vivrà negli anni 1902--1966 e non pubblicherà alcun lavoro.}. % lang-it

\index[persons]{Pedoe, Daniel}%
\index[persons]{Urquhart, Malcolm Livingston}%

Pompe napisze w $\Delta_{07}^2$ (bez przywołania nazwiska Urquharta) o możliwości dopisania okręgu do czworokąta. % lang-pl
Pompe scriverà in $\Delta_{07}^2$ (senza menzionare il nome di Urquhart) della possibilità di circoscrivere un cerchio a un quadrilatero. % lang-it

%